\chapter{Testing and Evaluation}

\section{Testing Strategy}

Testing is layered around business risk. Unit and service tests cover isolated
rules such as HMAC verification, merchant readiness, payment idempotency,
refund eligibility, expiration handling, webhook retry policy, and
reconciliation. Route tests cover FastAPI contracts and status codes. End-to-
end tests cover complete business journeys through controllers and services.
Frontend tests cover transforms, rendering, route behavior, type safety, and
build correctness.

\begin{longtable}{L{0.22\textwidth}L{0.34\textwidth}L{0.34\textwidth}}
\caption{Testing and evaluation matrix}\\
\toprule
\textbf{Category} & \textbf{Coverage} & \textbf{Evidence} \\
\midrule
\endfirsthead
\toprule
\textbf{Category} & \textbf{Coverage} & \textbf{Evidence} \\
\midrule
\endhead
Backend unit and service & Auth, readiness, payment, refund, expiration, callback, webhook, reconciliation, audit. & Python \texttt{unittest} suite under \code{backend/tests/}. \\
Backend route & Merchant APIs, provider callbacks, Ops APIs, Merchant Portal APIs, internal auth, internal users. & FastAPI controller and route tests. \\
Backend E2E & Onboarding to payment, refund, webhook, late callback reconciliation, and manual webhook retry. & \code{test_e2e_payment_refund_webhook.py}. \\
Merchant Portal scope & Auth, inactive user rejection, profile and credential secrecy, merchant-specific explorers, analytics ranges and zero buckets. & Merchant portal route and analytics tests. \\
Frontend Merchant Dashboard & Analytics transforms, chart render, route behavior, login accessibility, typecheck, build. & Vitest, TypeScript, and Vite build. \\
Frontend Ops Dashboard & Typecheck and production build. & TypeScript and Vite build. \\
Deployment configuration & Compose syntax and sandbox env requirements. & Docker Compose config validation. \\
Browser smoke & Login, dashboards, list and detail pages, analytics, mobile layout, and overflow checks. & Manual checklist in \code{docs/testing/README.md}. \\
\bottomrule
\end{longtable}

\section{Verification Commands}

\begin{lstlisting}[language=bash,caption={Backend verification}]
cd backend
python -m unittest discover tests -v
python -m unittest tests.test_e2e_payment_refund_webhook -v
\end{lstlisting}

\begin{lstlisting}[language=bash,caption={Frontend and compose verification}]
npm run merchant-dashboard:typecheck
npm run merchant-dashboard:test
npm run merchant-dashboard:build
npm run ops-dashboard:typecheck
npm run ops-dashboard:build
docker compose -f docker-compose.yml config --quiet
docker compose --env-file .env.sandbox.example -f docker-compose.sandbox.yml config --quiet
\end{lstlisting}

Historical pass counts and deploy verification snapshots are maintained in
\code{docs/history/completions/}. The report keeps the verification commands and test
coverage model in the foreground, while exact rollout-specific results stay in
the completion records.

\section{Manual Browser Smoke}

Manual smoke testing validates what automated tests cannot fully express:
layout, readable charts, responsive behavior, and practical operator workflows.
The smoke checklist includes:

\begin{itemize}
  \item Ops login, overview, merchant list and detail, lifecycle action,
  portal-user create or reset or deactivate, payments, refunds, webhooks,
  reconciliation, and audit;
  \item Merchant login, overview, analytics range switching, drill-down links,
  payments, refunds, webhooks, profile, credentials, password change, and
  logout;
  \item UI regressions such as horizontal overflow, clipped badges, unreadable
  charts, unstable status panels, and unreadable data tables.
\end{itemize}

\section{Evaluation Summary}

The current testing approach matches the system's main risks well:

\begin{itemize}
  \item merchant-facing business rules are covered at both service and route
  levels;
  \item asynchronous callback and webhook behavior is validated with focused
  service tests and end-to-end scenarios;
  \item dashboard correctness is supported by frontend build checks, targeted
  tests, and manual browser smoke;
  \item deployment wiring is validated separately through compose checks and
  sandbox deployment verification.
\end{itemize}

The remaining evaluation gap is typical for an MVP: browser smoke and internal
deployment verification still complement the automated suite because layout,
operational readability, and sandbox-specific runtime integration are not fully
captured by unit tests alone.
